\begin{figure}[p]
\centering
\includegraphics[width=\textwidth]{supplementary/curated/mechanistic_chain.pdf}
\caption{\textbf{Mechanistic chain from path gains to learning in regular-tree models.} MNIST panels \textbf{A--C} use two-stage $[4,4]$ directed trees; the historical noise-resilience panels \textbf{D,E} use two-stage $[3,3]$ additive trees. \textbf{A}, Conductance-stage path-gain field in shunting and normalized-additive models at five inhibitory synapses per branch; the inset shows paired-seed coefficients of variation (CV; standard deviation divided by mean). The schematic is drawn on the shared balanced tree of main Fig.~1B rather than the $[4,4]$ geometry: one root-to-leaf route carries the serial path gains \(\alpha_1\), \(\alpha_2\), \(\alpha_3\), distal sites are tinted by their illustrative path gain (darker, larger) in each model's colour, and \(\delta_0\) is the somatic error. This loading quantity omits branch derivatives and is distinct from the trained exact-error profiles in Supplementary Fig.~S8B,C. \textbf{B}, Stage-resolved cosine of matched-width scalar fallback versus exact MNIST errors; the width-matched somatic stage is excluded because its cosine is one by construction. \textbf{C}, MNIST post-training inhibition interventions: zero, sample shuffle, per-branch batch mean and a uniform matched mean. The dashed rule marks ten-class chance. These also change forward computation and do not isolate backward shunting. \textbf{D,E}, Feedback fidelity and learning on the additive noise task; solid and dashed lines denote matched-width scalar-fallback and exact transport. Bars show standard deviations across checkpoints in \textbf{B} and retained runs in \textbf{C}; shaded bands in \textbf{D,E} show standard deviations across retained runs (0.6--0.9 percentage points at doses 5--40 in \textbf{E}, and below 0.02 for the exact-path series in \textbf{D}, narrower than the marker); white points identify checkpoints in \textbf{B} and runs in \textbf{C}. Historical noise-task shunting aggregates with unresolved execution lineage are excluded and retained in Source Data only. The retained additive aggregates do not resolve their normalization setting. The state-matched focal analysis in main Fig.~8 provides the conductance-specific transport comparison. This chain is the figure support for Supplementary Section~S1; the exact-path field it feeds is drawn at the main-text type scale in main Fig.~1.}
\label{fig:si_mechanistic_chain}
\label{fig:inherited_mechanistic}
\end{figure}

\begin{figure}[p]
\centering
\includegraphics[width=\textwidth]{supplementary/curated/credit_validation.pdf}
\caption{\textbf{Local eligibility and exact transport reproduce the learning gradient.} \textbf{A}, Test accuracy on MNIST, Fashion-MNIST (F-MNIST) and figure-ground MNIST (FG-MNIST) for the shunting architecture trained by backpropagation (red-brown) or by the local rule (green), and for the normalized-additive architecture trained by the local rule (blue); bars are means $\pm$ 1 s.d. across five seeds, and the backpropagation intervals are narrower than the line weight. \textbf{B}, Transport, activation and additive-model implementation controls on MNIST; points are means $\pm$ 1 s.d. across five seeds, and the backpropagation interval is smaller than its marker. Colour names the architecture as in \textbf{A} (red-brown backpropagation, green shunting, blue additive, violet exact path transport); within a family the house marker (circle for shunting, square for additive) is the configuration drawn in \textbf{A} and a triangle is the ablated variant (identity activation without tanh reactivation; raw additive without normalization). The additive group draws the complete normalized/raw $\times$ gain-mode factorial, one gain mode per row. The MNIST backpropagation and normalized-additive conditions in \textbf{A} and \textbf{B} come from different run cohorts (the matched ceiling-refresh and local-rule tuning cohorts in \textbf{A}; the exact-transport and gain-normalization revision cohorts in \textbf{B}): the backpropagation means differ by 0.09 percentage points, and the normalized-additive bar in \textbf{A} lies within the range of the four normalized rows in \textbf{B}. \textbf{C}, Exact-transport-minus-backpropagation accuracy across three-/five-factor and local/backpropagated decoder choices; open circles are five paired seeds, diamonds are condition means and bars are 95\% paired-seed bootstrap intervals; dark violet denotes a backpropagated decoder and lighter violet a local decoder. Exact transport supplies an oracle compartment error. The local three-factor eligibility is sufficient when this error is supplied. Normalized- and raw-additive models are distinct controls; no noise-task shunting aggregate is used here.}
\label{fig:si_credit_validation}
\label{fig:inherited_competence}
\label{fig:inherited_rule_feedback}
\end{figure}

\begin{figure}[p]
\centering
\includegraphics[width=\textwidth]{supplementary/curated/utility_signal_noise.pdf}
\caption{\textbf{Restricted routes trade signal retention against admitted noise, and the one-step quantity is measurable at trained states.} \textbf{A}, Fixed-operator one-step utility bound for the update $M(\bm\mu+\bm\xi)$ (zero-mean noise $\bm\xi$): the boxed $U(M)$, the squared positive part of $\bm\mu^{\top}M\bm\mu$ over $2L\,\mathrm{E}\|M(\bm\mu+\bm\xi)\|^{2}$ ($L$, loss smoothness), lower-bounds the expected one-step loss decrease that \textbf{B--D} sweep. \textbf{B}, Subtree-minus-random spectral capture across route budget $K$ and covariance mixture $\rho$, printed per cell (red positive, blue negative, tint by magnitude); $K=16$ and $\rho=0$ are by-construction controls; only the endpoints are isospectral. \textbf{C}, Final quadratic loss against route resolution $D_{\rm r}$ for hierarchies $H_{\rm c}=1$--$4$, constructed so that fine-scale noise favours an interior resolution; $D_{\rm r}=4$ resolves every leaf ($M=I$); bands, 95\% bootstrap intervals, 50 paired seeds. \textbf{D}, Exact one-step projection boundary: retaining signal fraction $f_{\rm sig}$ and noise fraction $f_{\rm noise}$ changes normalized loss by $(f_{\rm noise}-f_{\rm sig})/2$, printed as in \textbf{B}; unequally spaced fractions form equal cells, so the dashed line is the cell-wise sign boundary, not the diagonal. \textbf{E}, Candidate/exact gradient cosine and norm-matched one-step loss decrease at the same 120 valid trained checkpoints (points; two off-scale strict-scalar values are printed at the floor); exact path is 1 by construction. Boxes, medians, quartiles and 1.5-interquartile-range whiskers; diamonds, means with 95\% checkpoint-bootstrap intervals (one held-out batch, relative step $10^{-5}$) narrower than the marker. \textbf{F}, Separate constructed positive control, eight reconstructed arbors: loss reduction after twenty projected steps versus the full-gradient sequence (green, morphology-selected routes; gray, random paths; blue, depth bins; pink, ancestry-shuffled routes); symbols, cell means with hierarchical 95\% bootstrap intervals. \textbf{G}, Illustrative rank--noise trade-off under isotropic noise, $2L$ times the optimized bound $q^2/(q+K\sigma^2)$ for $(K,q)=(1,0.8)$ (solid) and $(2,1)$ (dashed), crossing at $\sigma^2=4/7$ (shaded between); analytic, no data, not an endpoint-selection prediction. Imposed alignment and oracle projections explain the conditional relationship in \textbf{E,F}, which measures neither biological credit nor a long-horizon selector; the bound is local, and endpoint selection fails prospectively (Fig.~\ref{fig:si_original_selector}).}
\label{fig:si_utility_signal_noise}
\label{fig:supp_alignment}
\label{fig:supp_utility}
\label{fig:si_checkpoint_geometry}
\end{figure}

\begin{figure}[p]
\centering
\includegraphics[width=\textwidth]{supplementary/curated/reliability_gain.pdf}
\caption{\textbf{Fixed and estimated reliability gains have bounded learning benefits.} All panels use positive-rate synthetic tasks with eight parallel, nonserial branch blocks. \textbf{A}, A supplied compensating current preserves branch voltage while a positive shunt reduces input resistance and eligibility. \textbf{B}, Paired control-minus-aligned final-loss contrasts at maximal heterogeneity for the best global, shuffled, anti-aligned and unshunted controls; fixed aligned shunting has no reliable final advantage over unshunted noisy learning. \textbf{C,D}, Adaptive-rule final loss and paired control-minus-adaptive contrasts. Adaptive local gain beats global and shuffled attenuation but loses to no shunt and the fixed oracle. In \textbf{B,D} the dashed line is the aligned or adaptive-local reference, translucent dots are the 50 individual paired seeds, and diamonds with bars are means and 95\% paired-seed bootstrap intervals; an interval narrower than its diamond is hidden by it. The point implementation given identical gains reproduces the conductance final loss exactly in all 50 seeds of each study (paired difference 0, 50/50 ties). Curves and shading in \textbf{C} are means and 95\% paired-seed bootstrap intervals across the same 50 seeds. The supplied compensating current in \textbf{A} and the paired gradient observations from which the adaptive rule in \textbf{C,D} estimates its gains are explicit information resources.}
\label{fig:si_reliability_gain}
\label{fig:supp_positive_reliability}
\label{fig:supp_adaptive_reliability}
\end{figure}

\begin{figure}[p]
\centering
\includegraphics[width=\textwidth]{supplementary/curated/same_span_conditioning.pdf}
\caption{\textbf{Identical route spans can learn differently through conditioning.} \textbf{A}, Haar, raw nested (redundant) and scaled nested (statically scaled) coordinates parameterize the same rank-eight span in sixteen abstract coefficients. \textbf{B}, Positive Gram eigenvalues; condition numbers are 1, 15 and 388.52. \textbf{C}, Population loss after eighty updates across effective sample size on a logarithmic axis; curves and bands are means and 95\% seed-bootstrap intervals across fifty seeds. The strip above the axis shows, for every seed, the predicted effective sample size at which the exact finite-time risks of the nested and Haar coordinates are equal (raw nested: median 38.5, interquartile range 31.1--63.7; scaled nested: median 8.2, interquartile range 6.3--12.6); large symbols and dotted guides mark the medians. \textbf{D}, Paired nested-minus-Haar contrasts; negative values favor nested coordinates. Points and bars are means and 95\% paired-seed bootstrap intervals across fifty seeds (n = 50 pairs); dashed curves are the exact finite-time risk. Slow modes reduce variance at low data and retain bias at high data. Gram-preconditioned controls agree to $1.15\times10^{-15}$. This is a coefficient-learning comparison within one span, not a simulation of forward dendritic morphology or a biological implementation of the preconditioner.}
\label{fig:si_same_span_conditioning}
\label{fig:supp_same_span_learning}
\end{figure}
